\section{La compresión es inteligencia: por qué el lenguaje se volvió el soporte de la inteligencia}

La sección anterior trató cómo entiende el modelo un gerente de producto; esta sección aborda uno de los primeros principios de este episodio: la compresión es inteligencia. Según Wang Guan (王冠), la compresión une datos que antes eran puntos discretos en una estructura continua, y esa continuidad se manifiesta hacia fuera como generalización, emergencia, alucinación o inteligencia. El lenguaje importa porque puede expresar un mundo suficientemente rico con un costo de entrenamiento relativamente bajo.

\begin{figure}[H]
\centering
\includegraphics[width=0.96\textwidth]{compression-intelligence.png}
\caption{La compresión es inteligencia: la compresión une tareas discretas en una capacidad generalizable. Diagrama conceptual propio, elaborado a partir de la conversación de 00:28:39--00:32:25.}
\end{figure}

\begin{knowledgebox}{Lectura de la figura: entre los datos y la inteligencia no hay magia, sino continuidad}
En la figura, los «datos» son al principio muestras discretas; la compresión hace que formen relaciones continuas, el lenguaje aporta un soporte de alta capacidad expresiva, la generalización se manifiesta como la capacidad de resolver tareas no vistas y, al final, las personas llaman a eso inteligencia. Esta cadena explica por qué un modelo puede combinar «traducción del chino» y «resumen en chino» para obtener una capacidad como «resumen en inglés», que no se entrenó de forma explícita.
\end{knowledgebox}

\subsection{Por qué la compresión produce generalización}

En el programa se da un ejemplo intuitivo: si el modelo se entrenó en traducción del chino al inglés y también en resumen en chino, pero no en resumen en inglés, la representación comprimida puede permitirle unir las tareas de «traducir» y «resumir» y, así, producir un resumen en inglés. Esta capacidad no surge de la nada: en la distribución de los datos, las distintas tareas forman una estructura continua y transferible.

\begin{warningbox}{No hay que malinterpretar «la compresión es inteligencia» como una explicación universal}
La compresión explica una parte del origen de la generalización, pero eso no significa que cualquier compresión produzca inteligencia útil. La calidad de los datos, el objetivo de entrenamiento, la arquitectura del modelo, la capacidad de cómputo, la evaluación y el context del producto influyen en la capacidad final.
\end{warningbox}

\subsection{Implicaciones de producto de «el lenguaje es inteligencia»}

Más adelante, Wang Guan actualizó «la compresión es inteligencia» a «el lenguaje es inteligencia». Esta afirmación no niega las imágenes, el video ni las acciones; subraya que el lenguaje, como soporte abstracto, puede expresar a bajo costo el mundo, las tareas, los métodos, las relaciones y las intenciones. Para el producto, esto significa que el lenguaje natural no es solo una entrada de la UI, sino también un DSL, una descripción de tareas, un criterio de evaluación y una forma de organizar el context.

\subsection{Resumen de la sección}

La idea de que la compresión es inteligencia es la base de todos los juicios posteriores de este episodio: los datos determinan los límites, el lenguaje comprime el mundo en expresiones combinables, y el gerente de producto debe diseñar tanto las capacidades que el modelo puede aprender como el context que el modelo puede aprovechar.
